% ============================================================
% 分片：s05a-spectral  对应原文 PDF 第 56--66 页（印刷页 56--66）
% 负责章节：§5 Spectral Covers, Hecke Operators, and Higher Genus；5.1--5.2（§5 标题起）
% 译者：trans-s05a   审校状态：初稿
% ============================================================

% ==== 原文 PDF 第 56 页（印刷页 56）====

\section{谱覆盖、Hecke 算子与高亏格}

\begin{keytakeaway}
\textbf{核心物理图像：}Hitchin 纤维化的纤维可借助谱曲线实现为谱覆盖上满足次数条件的线丛（或无挠层）的模空间；内窥局部系统在 B-模型一侧给出 $\mathcal{M}_H(SO_3)$ 的 $\mathbb{Z}_2$ 轨形奇点，在 A-模型一侧给出对偶 Hitchin 纤维的两个分支，镜像对称把两个 B-膜 $\mathcal{B}_\pm$ 与两个 A-膜 $\mathcal{A}_1,\mathcal{A}_2$ 对应起来。
\textbf{数学工具：}谱曲线方程 $z^2+\det\varphi=0$、谱覆盖与 Prym 簇、torsor、$O_2^*\subset SL_2$ 的结构群约化、Riemann--Hurwitz 公式、Tsen 定理、群 $Q=H^1(C,\mathbb{Z}_2)$ 在 Hitchin 纤维上的作用、逐纤维 T-对偶。
\textbf{主要结论：}亏格 1 时奇异 Hitchin 纤维是两条 $\mathbb{CP}^1$ 之并，其正规化 $D'$ 是 $C$ 的非分歧二重覆盖；高亏格时内窥谱曲线的正规化是亏格 $2g-1$ 的非分歧二重覆盖，$SO_3$ 的 Hitchin 纤维模去 $Q$ 后只剩一个分支，而 $\mathbb{Z}_2$ 轨形奇点出现在余维 $4g-4$ 处；转移函子由 Hitchin 纤维化的镜像对称与逐纤维 T-对偶给出。
\end{keytakeaway}

第 3 节的分析主要基于直接计算。利用谱曲线的技巧可以学到多得多的东西，这一技巧已在 4.1 节中简要介绍。在 5.1 节中，我们从这一观点重新考察亏格 1 的例子；在 5.2 节中，我们把类似的想法应用于亏格大于 1 的情形；在 5.3 节中，我们利用这些方法分析 Hecke 算子对几何 Langlands 纲领中出现的 A-膜的作用。

\subsection{重访亏格 1}

我们首先回到第 3 节的例子，它涉及由三次方程
\begin{equation}\tag{5.1}
y^2 = f(x),
\end{equation}
所描述的亏格 1 曲线 $C$，其中 $f(x) = (x-e_1)(x-e_2)(x-e_3)$。我们考虑在由 $x=y=\infty$ 定义的点 $p$ 处分歧的 $SL_2$ Higgs 丛 $(E,\varphi)$。分歧意味着：在 $p$ 附近，$\varphi$ 的本征值表现为 $\pm\sigma_0\,dx/2x$，从而 $\det\varphi \sim -\sigma_0^2 dx^2/4x^2 \sim -(\sigma_0^2/4)x(dx/y)^2$。设 $\sigma_0\neq 0$。则 $C$ 上除在 $p$ 附近具有这种性态之外全纯的二次微分，其一般形式为 $\det\varphi = -(\sigma_0^2/4)(x-a)(dx/y)^2$，其中 $a$ 为某个复常数。于是，在把因子 $(dx/y)\sigma_0/2$ 吸收进 $z$ 之后，谱曲线 $D$ 的方程 (4.1) 变为
\begin{equation}\tag{5.2}
z^2 = x - a.
\end{equation}

% [Tutor Note: 推导补足。此处的「吸收因子」并非随意：谱曲线方程 (4.1) 为 $z^2+\det\varphi=0$，而 $\det\varphi$ 是二次微分，故 $z$ 必须是 $K_C$ 的截面（1-形式）。(5.2) 之所以能写成 $z^2=x-a$ 这种普通坐标形式，是因为用 $(dx/y)\sigma_0/2$ 重新标度了 $z$：原方程给出 $z^2=(\sigma_0^2/4)(x-a)(dx/y)^2$，除以 $(dx/y)^2\sigma_0^2/4$ 即得 (5.2)。注意 $\sigma_0$ 是常数，此处 $(dx/y)\sigma_0/2$ 是乘积而非幂；文字层抽出的线性串容易被误读成指数 $(dx/y)^{\sigma_0/2}$，务必以页面图像为准。 来源：原书 p.56 上下文（A）。]

% ==== 原文 PDF 第 57 页（印刷页 57）====

因此，$D$ 是由方程 (5.1) 与 (5.2) 这一对公式所描述的 $x$-平面的覆盖。用第二个方程消去 $x$，谱曲线可由方程
\begin{equation}\tag{5.3}
y^2 = \prod_{i=1}^{3}(z^2+a-e_i)
\end{equation}
描述。

$D$ 是一条亏格 2 的光滑曲线，除非 $a$ 等于 $e_1$、$e_2$ 或 $e_3$；在后一情形，$D$ 退化为一条亏格 1 的曲线，其中两个点被认同。例如，若 $a=e_1$，令 $w=y/z$，它满足
\begin{equation}\tag{5.4}
w^2 = (z^2+e_1-e_2)(z^2+e_1-e_3).
\end{equation}
该方程描述一条亏格 1 的光滑曲线 $D'$。$D$ 由 $D'$ 认同 $z=0$、$w=\pm\sqrt{(e_1-e_2)(e_1-e_3)}$ 的两个点得到。这两个点（我们将称之为 $q'$ 与 $q''$）对应于 $D$ 中的同一个点 $q$，因为它们都由 $y=z=0$ 刻画。

% [Tutor Note: 推导补足。(5.3) 表明 $D$ 是 $z$-平面的二重覆盖：$y^2$ 是 $z^2$ 的三次多项式，故有 6 个分歧点，Riemann-Hurwitz 给出 $2g(D)-2 = 2(-2)+6$，即 $g(D)=2$。当 $a=e_i$ 时其中一对分歧点合并，$D$ 出现一个节点；其正规化 $D'$ 由 (5.4) 给出（$z$-平面的二重覆盖，只剩 4 个分歧点），Riemann-Hurwitz 给出 $g(D')=1$，而 $D$ 由 $D'$ 认同 $q',q''$ 得到，算术亏格仍为 2。这就是下文「亏格 1 的曲线认同两个点」的含义。 来源：原书 p.57 上下文（A）。]

$D'$ 是 $C$ 的非分歧二重覆盖。事实上，$D'$ 具有自由作用的对称 $\tau: w\to -w,\ z\to -z$。不变量是 $x=z^2+e_1$，以及 $w^2$（它可借助式 (5.4) 用 $x$ 表示）与 $y=zw$（它满足 $y^2=\prod_{i=1}^{3}(x-e_i)$，即 $C$ 的定义方程）。所以 $C$ 就是商 $D'/\{1,\tau\}$。

我们用 $p'$ 与 $p''$ 记 $D'$ 中 $z=\infty$ 的两个点；它们位于 $C$ 中无穷远点 $p$ 之上。$D'$ 的对称 $\tau$ 自由作用且为 2 阶，因此若把 $D'$ 视为椭圆曲线、比方说以 $p'$ 为原点，则 $\tau$ 就是按一个 2 阶元素所作的平移。这个元素正是 $p''$，因为 $\tau$ 交换 $p'$ 与 $p''$。若 $r$ 是 $C$ 中任一点，而 $r',r''$ 是 $D'$ 中位于 $r$ 之上的点，则 $\tau$ 交换 $r',r''$，故
\begin{equation}\tag{5.5}
r'' - r' = p'' - p'.
\end{equation}

除子 $p'+p''-q'-q''$ 是函数 $z$ 的除子，故其除子类平凡。结合式 (5.5)，这说明点 $q',q''$ 都是 2 阶的。

\subsubsection{Hitchin 纤维化的纤维}

若 $D$ 光滑，Hitchin 纤维化的纤维是亏格 1 的光滑曲线，即二重覆盖 $\pi: D\to C$ 的 Prym 簇。我们来讨论当 $D$ 奇异时（例如在 $a=e_1$ 处）会发生什么。

我们要找出 $D$ 上的线丛，或者更一般地无挠层，使它们下推到 $C$ 上成为行列式平凡的向量丛 $E\to C$。我们将用光滑曲线 $D'$ 上的数据来描述这些线丛与层。因此，先从用投影 $\rho: D'\to C$「下推」平凡线丛 $\mathcal{O}_{D'}$ 开始。我们利用商 $D'/\{1,\tau\}$ 是 $C$ 这一事实来计算下推。可以把 $\rho_*(\mathcal{O}_{D'})$ 分解为在 $\tau$ 下为偶或为奇的子层。偶的部分是 $\mathcal{O}_C$，奇的部分是局部自由层 $\mathcal{S}\to C$。$\mathcal{S}$ 由如下事实唯一确定：它非平凡，且 $\rho_*(\mathcal{S}^{-1})$ 在 $D'$ 上有整体截面。我们可以用除子 $-p+q$ 来描述 $\mathcal{S}$，因为它在 $D'$ 上的拉回是平凡除子 $-p'-p''+q'+q''$。于是，$\mathcal{O}_{D'}$ 并不具有 $\rho_*(\mathcal{O}_{D'})$ 的行列式平凡这一性质。它的行列式是 $\mathcal{O}(p)^{-1}\otimes\mathcal{O}(q)$。然而，设 $p^*$ 为 $p',p''$ 中任一点，$q^*$ 为 $q',q''$ 中任一点。令 $\mathcal{T}_0 = \mathcal{O}(p^*)\otimes\mathcal{O}(q^*)^{-1}$。则 $\det(\rho_*(\mathcal{T}_0)) = \mathcal{O}_C$。此处我们需要的只是 $p^*$ 下推到 $p$、$q^*$ 下推到 $q$；这说明 $\det(\rho_*(\mathcal{T}_0)) =$
% ==== 原文 PDF 第 58 页（印刷页 58）====
$\mathcal{O}(p)\otimes\mathcal{O}(q)^{-1}\otimes\det(\rho_*(\mathcal{O}_{D'})) = \mathcal{O}_C$。看起来我们有 4 种选择 $\mathcal{T}$，但同构意义下只有 2 种；由式 (5.5)，若在反转 $q^*$ 的选择的同时也反转 $p^*$ 的选择，$\mathcal{T}$ 不变。

% [Tutor Note: 推导补足。所谓「4 种选择」指 $p^*\in\{p',p''\}$ 与 $q^*\in\{q',q''\}$ 的四种组合。由 (5.5)（取 $r=q$）得 $q''-q'=p''-p'$，于是取 $(p',q')$ 与取 $(p'',q'')$ 所对应的线丛相差 $\mathcal{O}(p''-p')\otimes\mathcal{O}(q'-q'') = \mathcal{O}(p''-p')\otimes\mathcal{O}(-(p''-p')) = \mathcal{O}$，二者同构；而取 $(p',q'')$ 与取 $(p'',q')$ 所对应的线丛相差 $\mathcal{O}(p''-p')\otimes\mathcal{O}(q''-q') = \mathcal{O}(2(p''-p'))$，这是非平凡的 2 阶线丛。故四种选择只给出两个同构类，这正是「只有 2 种」且它们给出的行列式都为 $\mathcal{O}_C$ 的原因。 来源：原书 p.58 上下文（A）。]

于是我们找到了两个下推到 $SL_2$-丛 $E\to C$ 的线丛 $\mathcal{T}_0\to D'$。它们是仅有的两个。但若在奇异曲线 $D$ 上（而非其正规化 $D'$ 上）工作，则有更多选择。我们先把次数 0 的线丛 $\mathcal{T}_0$ 换成次数 1 的线丛 $\mathcal{T}_1(p^*) = \mathcal{O}(p^*)$。我们把 $\mathcal{T}_1(p^*)$ 视为 $D'$ 上一个在 $p^*$ 之外被平凡化的线丛。我们不能把 $\mathcal{T}_1(p^*)$ 解释为 $D$ 上（而非 $D'$ 上）的线丛，因为在其定义中，我们没有考虑 $D'$ 上两个点 $q'$ 与 $q''$ 被认同为 $D$ 中单个点 $q$ 这件事。为此，取 $\lambda\in\mathbb{C}^\times$，并如下定义 $D$ 上的线丛 $\mathcal{T}(\lambda;p^*)$：$\mathcal{T}(\lambda;p^*)$ 的截面是 $\mathcal{T}_1$ 的截面 $f$，满足 $f(q') = \lambda f(q'')$。这给出（由于 $p^*$ 有两种选择）$D$ 上两族复线丛，每族以 $\mathbb{C}^\times$ 为参数。

为紧化这些族，令 $\lambda = -v/u$，其中 $u$ 与 $v$ 将被理解为 $\mathbb{CP}^1$ 的齐次坐标，并把条件 $f(q')=\lambda f(q'')$ 换成
\begin{equation}\tag{5.6}
u f(q') + v f(q'') = 0.
\end{equation}
对任意 $u,v$，满足该条件的 $\mathcal{T}_1(p^*)$ 的截面层是 $D$ 上的无挠层 $\mathcal{R}(u,v;p^*)$，其下推到 $C$ 是秩二、行列式平凡的向量丛。若 $u,v\neq 0$，该层局部自由，且是线丛 $\mathcal{T}(-v/u;p^*)$ 的截面层。若 $u$ 或 $v$ 为零，我们得到 $D$ 上无挠但非局部自由的层。这个无挠层是线丛 $\mathcal{O}(p^*)\otimes\mathcal{O}(q^*)^{-1}$ 从 $D'$ 的拉回，其中当 $v=0$ 时 $q^*=q'$（此时式 (5.6) 化为 $f(q')=0$），当 $u=0$ 时 $q^*=q''$。

借助关系 (5.5)（取 $r=q$），带不同 $p^*$ 与 $q^*$ 选择的四个线丛 $\mathcal{O}(p^*)\otimes\mathcal{O}(q^*)^{-1}\to D'$ 按两种选择成对同构，恰有两个同构类。因此，$\mathcal{R}(0,1;p')$ 同构于 $\mathcal{R}(1,0;p'')$，交换 $p',p''$ 后同理。

这一构造给出 $D$ 上两族无挠层 $\mathcal{R}(u,v;p^*)$。每族以 $\mathbb{CP}^1$ 为参数，齐次坐标为 $u,v$。由上一段的注记，这两个 $\mathbb{CP}^1$ 交于两个点。$SL_2$ 的 Hitchin 纤维化的纤维就是这两个 $\mathbb{CP}^1$ 之并。这与我们在 3.4 节通过直接计算得到的图景一致。

\subsubsection{improper 分支}

我们同样可以用谱曲线来描述 $SL_2$ 模空间的 improper 分支的 Hitchin 纤维化，该分支在 3.8 节引入。这个分支参数化行列式 $\det E = \mathcal{O}(r)$ 的 Higgs 丛 $(E,\varphi)$，其中 $r$ 是 $C$ 中一点。这里我们将处理 3.8 节略去的一个问题：对 $r$ 的依赖。因此我们把该 improper 分支记作 $\mathcal{M}_H(SL_2^*;r)$。

该空间在差一个并不完全典范的同构的意义下与 $r$ 无关；给定另一点 $\tilde{r}$，取线丛 $\mathcal{N}$ 使其平方同构于 $\mathcal{O}(\tilde{r})\otimes\mathcal{O}(r)^{-1}$，则用 $\mathcal{N}$ 作张量积给出从 $\mathcal{M}_H(SL_2^*;r)$ 到 $\mathcal{M}_H(SL_2^*;\tilde{r})$ 的映射。$\mathcal{N}$ 的选择在模去用 2 阶线丛作张量积的意义下唯一，故 $\mathcal{M}_H(SL_2^*;r)$ 与 $\mathcal{M}_H(SL_2^*;\tilde{r})$ 的等同在模去 2 阶线丛群 $Q$ 的作用的意义下唯一。（因此，若除以 $Q$，$\mathcal{M}_H(SL_2^*;r)$ 自然地与 $r$ 无关；这给出具有非零第二 Stieffel--Whitney 类的 $SO_3$ Higgs 丛的模空间 $\mathcal{M}_H(SO_3;\theta)$，其定义无需选择 $r$。）
% [原文疑误: 原文两处均拼作 "Stieffel-Whitney"，通行拼写为 "Stiefel-Whitney"（Stiefel 与 Whitney）；此处照录原文拼写，不代为改正。 来源：原书 p.58 英文原文与译本对照（A）。]
若 $\tilde{r}$ 靠近 $r$，我们可以要求 $\mathcal{N}$ 靠近恒等元（在 $C$ 的 Picard 群中）来消解这一歧义，于是局部地有 $\mathcal{M}_H(SL_2^*;r)$ 与 $\mathcal{M}_H(SL_2^*;\tilde{r})$ 的自然等同。因此存在自然的单值作用，单值群恰为 $Q$。

% [Tutor Note: 推导补足。$\mathcal{N}$ 之所以存在，是因为 $C$ 的 Picard 群可除：$\mathcal{O}(\tilde{r})\otimes\mathcal{O}(r)^{-1}$ 次数为 0，属于 $\mathrm{Pic}^0(C)$，而紧 Riemann 面的 $\mathrm{Pic}^0$ 是复环面，故可取平方根。平方根在模去 2 阶线丛（即 $\mathrm{Pic}^0$ 的 2-挠子群，同构于 $H^1(C,\mathbb{Z}_2)$）的意义下唯一，这就是下文的单值群恰为 $Q$ 的来源，也说明 $\mathcal{M}_H(SO_3;\theta)$ 是「除掉 $Q$ 之后」才自然定义的。 来源：原书 p.58 上下文（A）。]

% ==== 原文 PDF 第 59 页（印刷页 59）====

无论 $\det E$ 如何，Higgs 丛 $(E,\varphi)$ 的谱曲线都由方程 $\det(z-\varphi) = 0$ 定义。因此，对 Hitchin 纤维化的 improper 分支，相应的谱曲线 $D$ 与 proper 分支的相同。差别只在于：现在 Hitchin 纤维化的纤维参数化线丛 $\mathcal{R}\to D$，或更一般地无挠层，使得 $\det(\pi_*(\mathcal{R})) = \mathcal{O}(r)$（而非 $\det(\pi_*(\mathcal{R})) = \mathcal{O}$）。与前面一样，若 $D$ 光滑，则 Hitchin 纤维化的纤维光滑；我们要考虑的是那些 $D$ 不光滑、但以亏格 1 的光滑曲线 $D'$ 为正规化的特殊纤维。

为构造这些特殊纤维，我们重复先前的构造，但不再从线丛 $\mathcal{O}(p^*)\to D'$ 出发，而从 $\mathcal{T}_1 = \mathcal{O}(p^*)\otimes\mathcal{O}(r^*)$ 出发，其中 $r^*$ 是 $r',r''$ 中任一点。$\mathcal{T}_1$ 似乎有四种选择，涉及 $p^*$ 与 $r^*$ 的选择；但因为在椭圆曲线 $D'$ 上有 $r''-r' = p''-p'$，同构意义下只有两种选择。正如对模空间的 proper 分支那样，我们把这两个 $D'$ 上的线丛中的任一个与 $D$ 上的一族无挠层 $\mathcal{T}(u,v;p^*,r^*)$ 相联系，其办法是施加式 (5.6)。每族以 $\mathbb{CP}^1$ 为参数，并且如前，$\mathcal{T}(0,1;p^*,r')$ 同构于 $\mathcal{T}(1,0;p^*,r'')$，反之亦然。于是这两个 $\mathbb{CP}^1$ 交于两个点。它们的并就是 Hitchin 纤维化的纤维。

现在可以考虑当 $r$ 在 $C$ 中变化时的单值。这些单值将交换 $r^*$ 的两种选择，从而交换 Hitchin 纤维的两个分支。这与如下事实一致：单值群是 2 阶线丛群 $Q$，而 $Q$ 的作用交换 Hitchin 纤维的两个分支，正如我们在 3.5 节所见。

\subsection{向高亏格的推广}

我们的下一个目标是把同样的方法应用于 $C$ 是亏格 $g>1$ 的光滑曲线的情形。我们将看到，结果类似。为简单起见，我们略去分歧。对 $g=1$，非分歧情形相当退化，但对 $g>1$ 则不然。

\subsubsection{谱曲线}

什么样的 $C$ 上 Higgs 丛会与内窥相关？我们考虑结构群约化到 $O_2$ 的内窥 $SO_3$ 局部系统，其中 $O_2$ 是形如
\begin{equation}\tag{5.7}
\begin{pmatrix}
* & * & 0\\
* & * & 0\\
0 & 0 & \pm 1
\end{pmatrix},
\end{equation}
的元素构成的子群，但不约化到真子群。对这种局部系统，自同构群等于 $\mathbb{Z}_2$。Hitchin 方程给出的局部系统与 Higgs 丛之间的对应与任何结构群约化相容。因此，结构群约化到 $O_2$ 的 $SO_3$ 局部系统对应于具有相同结构群的 Higgs 丛。那些不再进一步约化到真子群的，其自同构群等于 $\mathbb{Z}_2$。从现在起我们将只限于这些 Higgs 丛。

% ==== 原文 PDF 第 60 页（印刷页 60）====

若把这样的 Higgs 丛提升到 $SL_2$，结构群提升到我们称之为 $O_2^*$ 的群，即由 $SL_2$ 的对角元素
\begin{equation}\tag{5.8}
\begin{pmatrix}
* & 0\\
0 & *
\end{pmatrix}
\end{equation}
与元素
\begin{equation}\tag{5.9}
\begin{pmatrix}
0 & 1\\
-1 & 0
\end{pmatrix}
\end{equation}
生成的 $O_2$ 的二重覆盖。

$SL_2$ 谱曲线 $D$ 由方程
\begin{equation}\tag{5.10}
z^2 + \det\varphi = 0
\end{equation}
定义，其中 $\det\varphi$ 是 $C$ 上的二次微分。对亏格 $g$ 的 $C$，二次微分有 $4g-4$ 个零点。$D$ 光滑当且仅当 $\det\varphi$ 的零点互不相同；此时 $D$ 的亏格为 $4g-3$。

当结构群约化到 $O_2^*$ 时会发生什么？$O_2^*$ 的 Lie 代数就是无迹对角矩阵构成的阿贝尔代数。因此，若 $(E,\varphi)$ 是结构群约化到 $O_2^*$ 的 $SL_2$ Higgs 丛，则 $\varphi$ 局部地形如
\begin{equation}\tag{5.11}
\varphi = \begin{pmatrix}
\omega & 0\\
0 & -\omega
\end{pmatrix},
\end{equation}
其中 $\omega$ 是 $C$ 上的全纯微分。这导致 $\det\varphi = -\omega^2$，因此 $\det\varphi$ 的零点都是二重零点。在例如 $x=0$ 处的二重零点附近（$x$ 是 $C$ 上的局部参数），$D$ 的方程看起来像 $z^2-x^2=0$；点 $z=x=0$ 是一个二重点。每个二重点使 $D$ 的正规化的亏格减少 1。对 $O_2^*$ Higgs 丛，零点全是二重零点，故有 $2g-2$ 个二重点，而 $D$ 的正规化是亏格 $2g-1$ 的曲线 $D'$。

% [Tutor Note: 推导补足。亏格计数如下：$D'$ 是 $C$ 的非分歧二重覆盖，Riemann-Hurwitz 给出 $2g(D')-2 = 2(2g-2)$，故 $g(D')=2g-1$；而 $D$ 由 $D'$ 认同 $\omega$ 的每个零点上方的两个点（共 $2g-2$ 对）得到，每认同一对点使算术亏格增加 1，于是 $g(D) = (2g-1)+(2g-2) = 4g-3$，与 $D$ 光滑时的 $4g-3$ 相符。注意这里的二重点都来自 $\omega$ 的二重零点，而不是任意节点，这正是内窥情形谱曲线的特殊之处。 来源：原书 p.60 上下文（A）。]

$\det\varphi$ 的零点都是二重零点这一事实，并不意味着整体上 $\det\varphi = -\omega^2$，其中 $\omega$ 是典范丛 $K$ 的全纯截面。它实际意味着 $\det\varphi = -\omega^2$，其中 $\omega$ 是 $K\otimes\mathcal{V}$ 的全纯截面，$\mathcal{V}\to C$ 是某个 2 阶线丛。我们关心的是 $\mathcal{V}$ 非平凡的情形。（$\mathcal{V}$ 平凡的情形与结构群约化到 $GL_1$ 而非 $O_2^*$ 的 Higgs 丛相关。）与 $\mathcal{V}$ 的选择相联系的是非分歧二重覆盖 $\pi': D'\to C$。这是亏格 $2g-1$ 的曲线，也是 $D$ 的正规化。设 $\tau: D'\to D'$ 是与 $\pi'$ 交换的覆盖映射。则 $C = D'/\{1,\tau\}$。元素 $\omega\in H^0(C,K\otimes\mathcal{V})$ 在 $2g-2$ 个点 $p_1,\dots,p_{2g-2}$ 处为零。在每个这样的点 $p_i$ 之上有 $D'$ 中的两个点 $p'_i,p''_i$，它们被 $\tau$ 交换，但 $D$ 中只有一个点 $\hat{p}_i$。$D$ 由 $D'$ 把成对的点 $p'_i$ 与 $p''_i$ 粘合而得。

刚才所描述的是 5.1 节结果对带一个分歧点的亏格 1 曲线 $C$ 的自然类比。亏格 1 的曲线恰有三个非平凡非分歧二重覆盖；它们就是三条奇异谱曲线 $D$ 的正规化 $D'$。

\subsubsection{Prym 簇}

$SL_2$ 的 Hitchin 纤维 $\boldsymbol{F}$ 的一个开稠子集由满足 $\det\pi_*(\mathcal{L}) = \mathcal{O}$ 的线丛 $\mathcal{L}\to D$ 构成。完整的 Hitchin 纤维除了这些线丛之外，还参数化某些无挠层。

% ==== 原文 PDF 第 61 页（印刷页 61）====

给定这样的线丛 $\mathcal{L}$，设 $\mathcal{L}'$ 为其到 $D'$ 的拉回。则 $\det\pi'_*(\mathcal{L}') = \mathcal{O}(p_1+\cdots+p_{2g-2})$。（当 $\mathcal{L}$ 提升到 $D'$ 时，我们放弃「截面在 $p_i$ 上方的点 $p'_i$ 与 $p''_i$ 处必须取相同值」这一要求，这导致所述结果。）这样的线丛构成的空间是\footnote{这句话的意思是：若 $\mathcal{L}'$ 与 $\mathcal{L}''$ 是两个线丛，且 $\det\pi'_*(\mathcal{L}') = \det\pi'_*(\mathcal{L}'') = \mathcal{O}(p_1+\cdots+p_{2g-2})$，则 $\mathcal{L}'' = \mathcal{L}'\otimes\mathcal{N}$，其中 $\mathcal{N}$ 唯一且满足 $\det\pi'_*(\mathcal{N}) = \mathcal{O}$。} 线丛群 $\mathcal{N}\to D'$（满足 $\det\pi'_*(\mathcal{N}) = \mathcal{O}$）的一个「torsor」。由于 $\pi': D'\to C$ 非分歧，$\det\pi'_*(\mathcal{N})$ 与 $\mathrm{Nm}(\mathcal{N})$ 相同，后者是 $\mathcal{N}$ 的\footnote{对任何曲线映射 $\pi: D\to C$，范数是从 $D$ 上的线丛到 $C$ 上的线丛的映射，定义如下。对整数 $n_i$ 与点 $p_i\in D$，$\mathcal{N} = \mathcal{O}(\sum_i n_i p_i)\to D$ 的范数定义为 $\mathrm{Nm}(\mathcal{N}) = \mathcal{O}(\sum_i n_i \pi(p_i))$。}「范数」。平凡范数的线丛群称为 $\pi': D'\to C$ 的 Prym 簇。（有时「Prym 簇」一词被用来指这个群的连通分支。）我们把 Prym 记作 $\boldsymbol{P}$，把满足 $\det\pi'_*(\mathcal{L}') = \mathcal{O}(p_1+\cdots+p_{2g-2})$ 的线丛 $\mathcal{L}'\to D'$ 所构成的 torsor 记作 $\boldsymbol{T}$。

我们可以容易地构造 $D'$ 上平凡范数线丛的一大族。取任意点 $s_i\in C$，$i=1,\dots,k$（允许某些点重合），把 $D'$ 中位于 $s_i$ 之上的点记作 $s'_i$ 与 $s''_i$（哪一个记作哪一个可任意选取），并如下定义线丛 $\mathcal{N}$：
\begin{equation}\tag{5.12}
\mathcal{N} = \bigotimes_{i=1}^{k}\mathcal{O}(s'_i - s''_i).
\end{equation}
反之，可以证明 $\boldsymbol{P}$ 的每一点都可由这种形式的线丛表示。为看清这一点（我们要感谢 T. Pantev 向我们展示这一论证），设 $E$ 是 $D'$ 上的除子，使得 $\mathcal{N} = \mathcal{O}(E)$。则 $\mathrm{Nm}(\mathcal{N}) = \mathcal{O}_C(\pi'(E))$。若 $\mathcal{N}\in\boldsymbol{P}$，则 $\pi'(E) = (f)$，即 $C$ 上有理函数 $f$ 的除子。但由 Tsen 定理，从 $D'$ 上有理函数域到 $C$ 上有理函数域的范数映射是满的\footnote{有理函数 $g$ 在 $D'$ 上的范数按定义为 $g\cdot\tau(g)$ 的乘积，其中 $\tau$ 是对应于覆盖 $D'\to C$ 的对合。}。于是存在 $D'$ 上的有理函数 $g$，其范数为 $f$。则 $\pi'((g)) = (f)$。令 $E' = D-(g)$。则 $\pi'(E') = 0$，因此 $E'$ 具有 $\sum_{i=1}^{k}(s'_i-s''_i)$ 的形式。这样就得到式 (5.12)。

% [原文疑误: 此处原文作 $E' = D-(g)$，应为 $E' = E-(g)$：由 $\pi'(E) = (f)$ 与 $\pi'((g)) = (f)$ 得 $\pi'(E-(g)) = 0$，故 $E'$ 应是 $E$ 减去主除子 $(g)$；若取 $D-(g)$ 则无法得到 $\pi'(E')=0$（$D$ 是曲线而非除子）。照录原文，不做改动。 来源：原书 p.61 英文原文与译本对照（A）。]

% [Tutor Note: 推导补足。Tsen 定理在此的作用是：代数闭域上曲线的函数域是 $C_1$ 域，而 $C_1$ 域的有限扩张上范数映射 $(F')^\times\to F^\times$ 满。等价的表述是 $H^1(\mathbb{Z}_2,(F')^\times)=0$，即 Hilbert 90 的推广。这一步正是「$\boldsymbol{P}$ 的每一点都可用 (5.12) 表示」的全部实质：先在除子层面把范数平凡化为主除子，再用范数映射的满性把主除子提升为 $D'$ 上的主除子，最后移去后者即得形如 $\sum(s'_i-s''_i)$ 的除子。 来源：原书 p.61 上下文（A）。]

现在我们来找出 Prym 的连通分支。显然，对固定 $k$ 能构造出的 $\boldsymbol{P}$ 的那一部分是连通的，因为 $C$ 本身连通且点 $s_i$ 可以自由移动。相差一次交换 $s'_i\leftrightarrow s''_i$ 的线丛位于 $\boldsymbol{P}$ 的同一连通分支，因为这两点在 $s_i$ 的单值作用下被交换。相差把 $k$ 改变 2 的倍数的线丛也位于同一连通分支；若取两个 $s_i$ 相等、把点适当标记为 $s'_i$ 与 $s''_i$，并利用恒等式 $(s'_i-s''_i)+(s''_i-s'_i) = 0$，$k$ 就减少 2。

另一方面，$\boldsymbol{P}$ 实际上有两个连通分支。这在 [Mum] 中被证明，也可由纯拓扑的论证得出（参见对式 (5.15) 的讨论）。\footnote{另一个证明见 [N1] 第 11 节。}上一段的论证表明，唯一可能的不变量是 $k$ 模 2 的值。因此必定是：$\boldsymbol{P}$ 中与恒等元相连的分支由偶数 $k$ 刻画，而不连通分支由奇数 $k$ 刻画。

% [Tutor Note: 易错预警。(5.12) 中 $k$ 的奇偶是 Prym 分支的唯一不变量。$k=0$ 给出平凡线丛，故偶 $k$ 落在含恒等元的分支；$k=1$ 给出 $\mathcal{O}(s'-s'')$，落在另一分支。注意 $k$ 的奇偶只是分支的标签（即商 $\boldsymbol{P}/\boldsymbol{P}^0\cong\mathbb{Z}_2$ 中的类），不要误以为 $\mathcal{O}(s'-s'')$ 是 2 阶线丛：一般地 $2(s'-s'')$ 并不主，故它通常不是 2-挠元。这一点与 5.1 节亏格 1 的情形形成对照，那里 $q'-q''$ 恰是 2 阶的（见 (5.5) 后的讨论）。 来源：原书 p.61 上下文（A）。]

% ==== 原文 PDF 第 62 页（印刷页 62）====

torsor $\boldsymbol{T}$ 非典范地同构于 $\boldsymbol{P}$，故它同样有两个分支。我们按 5.1.1 节的方式，从 $\boldsymbol{T}$ 构造 Hitchin 纤维化的纤维 $\boldsymbol{F}$。设 $\mathcal{L}'$ 是代表 $\boldsymbol{T}$ 中任一点的线丛。对 $i=1,\dots,2g-2$，取一对齐次坐标 $(u_i,v_i)$，并如下定义线丛 $\mathcal{L}'(u_1,v_1;\dots;u_{2g-2},v_{2g-2})\to D$：其截面是 $D'$ 上线丛 $\mathcal{L}'$ 的截面 $f$，满足
\begin{equation}\tag{5.13}
u_i f(p'_i) + v_i f(p''_i) = 0,\qquad i = 1,\dots,2g-2.
\end{equation}
若对所有 $i$，$u_i$ 与 $v_i$ 都非零，则这一构造给出一族 $D$ 上的线丛；若再让 $\mathcal{L}'$ 在 $\boldsymbol{T}$ 中变动，它们就表示 Hitchin 纤维 $\boldsymbol{F}$ 的一个 Zariski 开集。这个开集有两个连通分支 $\widetilde{\boldsymbol{F}}_1$ 与 $\widetilde{\boldsymbol{F}}_2$，因为 $\mathcal{L}'$ 可位于 $\boldsymbol{T}$ 的任一分之中。它们同构于 $\boldsymbol{T}$ 的两个分支上的 $(\mathbb{C}^\times)^{2g-2}$-丛。

如同 5.1.1 节，为得到完整的 Hitchin 纤维，我们必须通过纳入允许某个 $i$ 的 $u_i$ 或 $v_i$ 为零所得的无挠层来紧化。紧化后的纤维有两个不可约分支，我们称之为 $\boldsymbol{F}_1$ 与 $\boldsymbol{F}_2$。它们同构于 $\boldsymbol{T}$ 的两个分支上的 $(\mathbb{CP}^1)^{2g-2}$-丛。然而，正如亏格 1 的例子那样，$\boldsymbol{F}_1$ 与 $\boldsymbol{F}_2$ 在 $u_i$ 或 $v_i$ 为零的除子上相交。原因是：从线丛 $\mathcal{L}'\to D'$ 出发、在上一段的构造中取 $u_i=0$，等价于从另一个线丛 $\mathcal{L}'' = \mathcal{L}'\otimes\mathcal{O}(p'_i-p''_i)$ 出发并取 $v_i=0$。但操作 $\mathcal{L}'\to\mathcal{L}''$ 交换 $\boldsymbol{T}$ 的两个连通分支。

% [Tutor Note: 推导补足。相交的原因是粘合条件的齐次性：对同一个 $f$，「$u_i=0$ 且 $v_i\neq0$」迫使 $f(p''_i)=0$，而「$v_i=0$ 且 $u_i\neq0$」迫使 $f(p'_i)=0$；把 $\mathcal{L}'$ 换成 $\mathcal{L}'' = \mathcal{L}'\otimes\mathcal{O}(p'_i-p''_i)$ 恰好把这两种消失条件互换，于是两族在 $u_i=0$ 或 $v_i=0$ 的除子上重合。这也解释了为什么紧化后的纤维是可约的：两个分支沿 $2g-2$ 个除子粘在一起，形成亏格 1 情形中两条 $\mathbb{CP}^1$ 交于两点的推广。 来源：原书 p.62 上下文（A）。]

\subsubsection{improper 分支}

我们还想理解 $\det E = \mathcal{O}(r)$（$r$ 是 $C$ 中指定的一点）的 Higgs 丛 $(E,\varphi)$。此时 Hitchin 纤维可用同样的方法分析。我们只需从同一 Prym 簇 $\boldsymbol{P}$ 的另一个 torsor $\boldsymbol{T}(r)$ 出发。$\boldsymbol{T}(r)$ 参数化满足 $\det\pi'_*(\mathcal{L}') = \mathcal{O}(p_1+\cdots+p_{2g-2}+r)$ 的线丛 $\mathcal{L}'\to D'$。$\boldsymbol{T}(r)$ 同样有两个连通分支，它们被与 $\mathcal{O}(s'-s'')$（$s\in C$ 任意）作张量积所交换。相应地，Hitchin 纤维化的纤维有两个分支，沿参数化非局部自由无挠层的除子相交。

设 $\boldsymbol{F}$ 与 $\boldsymbol{F}^*(r)$ 分别是 $\det E = \mathcal{O}$ 与 $\det E = \mathcal{O}(r)$ 的 Higgs 丛 $(E,\varphi)$ 的 Hitchin 纤维。则 $\boldsymbol{F}$ 与 $\boldsymbol{F}^*(r)$ 非典范同构。为作出一个同构，只需取 $D'$ 中位于 $r$ 之上的点 $r'$。然后，对 $\mathcal{L}'\in\boldsymbol{T}$，通过 $\mathcal{L}'_{r'} = \mathcal{L}'\otimes\mathcal{O}(r')$ 定义 $\boldsymbol{T}(r)$ 中的线丛 $\mathcal{L}'_{r'}$。我们用 $\mathcal{L}'(u_i,v_i)\to\mathcal{L}'_{r'}(u_i,v_i)$ 把 $\boldsymbol{F}$ 映到 $\boldsymbol{F}^*(r)$。这是 $\boldsymbol{F}$ 与 $\boldsymbol{F}^*(r)$ 之间的同构，但由于它依赖于 $r'$ 的选择，并不完全典范。

现在考虑 $r$ 中交换 $r'$ 与 $r''$ 的单值的效果。其效果是把 $\mathcal{L}'_{r'}$ 映到 $\mathcal{L}'_{r''} = \mathcal{L}'_{r'}\otimes\mathcal{O}(r''-r')$。这一操作交换 $\boldsymbol{F}^*(r)$ 的两个分支。于是，$r$ 中的单值交换 $\boldsymbol{F}^*(r)$ 的两个分支，正如我们在 5.1.2 节末尾的亏格 1 例子中所见。

\subsubsection{拓扑观点}

这里我们将从拓扑观点解释如下事实：对非分歧（但连通）的二重覆盖 $D'\to C$，其 Prym 簇 $\boldsymbol{P}$ 有两个分支。

这个 Prym 是某些 $O_2^*$ 局部系统的 Hitchin 纤维化的纤维。我们回忆，$O_2^*$ 是 $SL_2$ 中由对角矩阵与元素
% ==== 原文 PDF 第 63 页（印刷页 63）====
\begin{equation}\tag{5.14}
\begin{pmatrix}
0 & 1\\
-1 & 0
\end{pmatrix}
\end{equation}
生成的子群。在二重覆盖 $SL_2\to SO_3$ 下，$O_2^*$ 投影到 $O_2\subset SO_3$。拓扑上，$O_2^*$ 有两个分支；含恒等元的分支由对角元素构成，另一分支由 $O_2^*$ 中非对角的元素构成。

Prym 有两个分支这一陈述，等价于：即使选定了非分歧二重覆盖 $D'\to C$，$\mathcal{M}_H(O_2^*)$ 的相应分支实际上仍有两个分支。（对给定 $D'$，Hitchin 纤维化的底连通，故 $\mathcal{M}_H(O_2^*)$ 的分支就是 Prym 的分支。）利用 $O_2^*$ Higgs 丛与 $O_2^*$ 局部系统的关系，确定 $\mathcal{M}_H(O_2^*)$ 的分支的问题等价于关于 $O_2^*$ 局部系统的类似问题，并可用拓扑方式回答。

取 $C$ 的基本群的适当生成元，并把 $O_2^*$ 局部系统的单值记作 $A_i,B_j$，$i,j=1,\dots,g$，我们有
\begin{equation}\tag{5.15}
[A_1,B_1][A_2,B_2]\cdots[A_g,B_g] = 1,
\end{equation}
其中 $[A,B] = ABA^{-1}B^{-1}$。我们这样来指定 $\mathcal{V}$：例如说 $B_g$ 位于 $O_2^*$ 的非连通分支，而其余所有 $A_i$ 与 $B_j$ 位于连通分支。由于连通分支是阿贝尔的，式 (5.15) 化为 $A_g B_g A_g^{-1}B_g^{-1} = 1$；由此推出（对位于非连通分支的 $B_g$）$A_g$ 是 $O_2^*$ 的两个中心元素之一。$\boldsymbol{P}$ 的两个分支由 $A_g$ 的选择来区分。

% [Tutor Note: 推导补足。作者省略的关键一步是：$O_2^*$ 的连通分支是对角环面，而用非对角元素共轭会把对角元素变为其逆，即 $B_g A_g B_g^{-1} = A_g^{-1}$。由 $A_gB_gA_g^{-1}=B_g$ 左乘 $A_g^{-1}$ 得 $B_gA_g^{-1}=A_g^{-1}B_g$，即 $B_g$ 与 $A_g^{-1}$（从而与 $A_g$）交换；与「共轭取逆」联立得 $A_g = A_g^{-1}$，即 $A_g = \pm 1$，正是两个中心元素。因此 $\boldsymbol{P}$ 的两个分支恰好由 $A_g$ 的这个符号区分。 来源：原书 p.63 上下文（A）。]

\subsubsection{\texorpdfstring{$SO_3$}{SO3} 的类比}

为理解内窥，我们必须考虑 $SO_3$（而非 $SL_2$）的 Hitchin 纤维化的纤维。正如式 (3.20) 那样，$SO_3$ Higgs 丛的模空间由 $SL_2$ Higgs 丛的模空间除以群 $Q = H^1(C,\mathbb{Z}_2)$（2 阶线丛群）得到。$Q$ 作用在 Higgs 丛 $(E,\varphi)$ 上为 $E\to E\otimes\mathcal{R}$，其中 $\mathcal{R}$ 是 2 阶线丛。

这一操作不影响 $\det\varphi$，故 $Q$ 的作用与 Hitchin 纤维化交换。$Q$ 在 Hitchin 纤维化的一个纤维上的作用容易描述：若 $E = \pi_*(\mathcal{L})$，其中 $\mathcal{L}$ 是谱覆盖上的线丛，则 $E\otimes\mathcal{R} = \pi_*(\mathcal{L}\otimes\pi^*(\mathcal{R}))$。于是 $Q$ 在 Hitchin 纤维化的纤维上的作用为 $\mathcal{L}\to\mathcal{L}\otimes\pi^*(\mathcal{R})$。

在谱曲线 $D$ 光滑的情形，Hitchin 纤维化的纤维 $\boldsymbol{F}$ 是复环面，而操作 $\mathcal{L}\to\mathcal{L}\otimes\pi^*(\mathcal{R})$ 是这个环面上的平移。它无不动点地作用。现在来看与内窥相关的情形，即 $D$ 的正规化是非分歧二重覆盖 $D'\to C$ 的情形。如前所见，$\boldsymbol{F}$ 可用 torsor $\boldsymbol{T}$ 方便地描述，\footnote{$Q$ 在 improper 纤维 $\boldsymbol{F}^*(r)$ 上的作用可类似地考虑。}后者参数化满足 $\det\pi_*(\mathcal{L}') = \mathcal{O}(p_1+\cdots+p_{2g-2})$ 的线丛 $\mathcal{L}'\to D'$。$Q$ 通过 $\mathcal{L}'\to\mathcal{L}'\otimes\mathcal{R}$ 作用在这个 torsor 上。$Q$ 也作用在式 (5.13) 的粘合数据上，我们稍后马上讨论。
% [存疑: 此处原文作 $\det\pi_*(\mathcal{L}')$（$\pi$ 无撇号），而 5.2.2 节同一条件写作 $\det\pi'_*(\mathcal{L}')$；因 $\mathcal{L}'$ 定义在 $D'$ 上，疑为笔误（或作者的简写）。照录原文。 来源：原书 p.63 英文原文；存疑待核（A）。]

关于这一内窥情形的第一个基本事实是：[Mum] $Q$ 的作用交换 Prym 簇 $\boldsymbol{P}$ 的两个分支，从而交换 torsor $\boldsymbol{T}$ 的两个分支。（当然，$Q$ 的指数 2 子群把给定分支映到自身。）这一事实用式 (5.15) 有简单的拓扑解释。$Q$ 通过对所有 $A_i$ 与 $B_j$ 作独立的变号来作用；
% ==== 原文 PDF 第 64 页（印刷页 64）====
$\boldsymbol{P}$ 的两个分支被 $Q$ 中任何改变 $A_g$ 的符号的元素所交换。

于是在除以 $Q$ 之后，$SO_3$ 的 Hitchin 纤维只有一个分支，这与 $SL_2$ 的情形相反。这一点从 3.6 节已熟悉。

第二点同样从 3.6 节已熟悉：$SO_3$ Higgs 丛的模空间具有奇点，其来源是 $Q$ 的作用并不完全自由。这如何发生？若 $\mathcal{R}\to C$ 是 2 阶线丛，且其到 $D'$ 的拉回非平凡，则操作 $\mathcal{L}\to\mathcal{L}\otimes\mathcal{R}$ 自由地作用在 $\boldsymbol{T}$ 上，从而自由地作用在 $\boldsymbol{F}$ 上。

然而，$\pi': D'\to C$ 是与某 2 阶线丛 $\mathcal{V}\to C$ 相联系的非分歧二重覆盖，而同义地，$\mathcal{V}$ 到 $D'$ 的拉回 $(\pi')^*(\mathcal{V})$ 平凡。于是 $Q$ 中对应于 $\mathcal{V}$ 的元素平凡地作用在 $\boldsymbol{T}$ 上。但它非平凡地作用在式 (5.13) 的粘合条件上。$(\pi')^*(\mathcal{V})$ 平凡意味着该线丛有处处非零的整体截面 $w$。这个截面在覆盖映射 $\tau: D'\to D'$ 下是奇的而非偶的。（否则 $w$ 会下降为 $\mathcal{V}\to C$ 的截面，与 $\mathcal{V}$ 非平凡矛盾。）$\mathcal{V}$ 的作用把式 (5.13) 的粘合条件改成 $u_i\tilde{f}(p'_i) + v_i\tilde{f}(p''_i) = 0$，其中 $\tilde{f} = f w$。由于 $w$ 在交换所有 $i$ 的 $p'_i$ 与 $p''_i$ 的覆盖映射下为奇，其效果是把粘合数据变换为
\begin{equation}\tag{5.16}
(u_i,v_i)\to (u_i,-v_i),\qquad i=1,\dots,2g-2.
\end{equation}
记住 $u_i,v_i$ 是一份 $\mathbb{CP}^1$ 的齐次坐标，$\mathcal{V}$ 的不动点条件是：对所有 $i$，$u_i$ 或 $v_i$ 必须为零。这是 $2g-2$ 个条件。

我们还调整了 $2g-2$ 个参数，使谱曲线 $D$ 以二重覆盖 $D'\to C$ 为其正规化。作为检验，同时也为了确认这类不动点只在我们所假定的那类谱曲线上出现，我们注意到：$Q$ 作用保持 Higgs 丛模空间的复辛结构，而 Hitchin 纤维是拉格朗日的，这蕴含着——为使某点成为不动点而必须在 Hitchin 纤维化的纤维上调整的参数个数，等于必须在底上调整的参数个数。

% [Tutor Note: 推导补足。这里的维数论证是超凯勒几何中的标准事实：Hitchin 纤维化是复辛流形之间的拉格朗日纤维化，在不动点处辛形式把纤维方向的法空间与底方向的切空间配对，于是「纤维上需调整的参数个数 = 底上需调整的参数个数」。具体地，(5.16) 中每个 $(\mathbb{CP}^1)$ 因子有两个不动点 $u_i=0$ 与 $v_i=0$，共 $2g-2$ 个条件；底上则要求 $\det\varphi$ 的 $2g-2$ 个零点都成为二重零点（即 $\det\varphi$ 落在 $\omega^2$ 的像中）。两者相加给出总余维 $2(2g-2)=4g-4$。 来源：原书 p.64 上下文（A）。]

总之，我们所找到的 $\mathbb{Z}_2$ 不动点出现在余维 $4g-4$ 处。局部结构是 $\mathbb{C}^{4g-4}/\mathbb{Z}_2$。这些奇点不是 $A_1$ 奇点，因为它们出现在大于 2 的复余维上。但 2.2 节的一般性论述仍然适用。在任意余维的 $\mathbb{Z}_2$ 轨形奇点处，有两个不等价的 B-膜 $\mathcal{B}_+$ 与 $\mathcal{B}_-$；我们的基本假设是，\emph{镜像对称把这一事实映为如下事实：$SL_2$ 的相应 Hitchin 纤维有两个分支 $\boldsymbol{F}_1$ 与 $\boldsymbol{F}_2$}。于是，两个 B-膜 $\mathcal{B}_+$ 与 $\mathcal{B}_-$ 给出支撑在这两个分支上的两个 A-膜 $\mathcal{A}_1$ 与 $\mathcal{A}_2$。其中哪一个对应于 $\mathcal{B}_+$、哪一个对应于 $\mathcal{B}_-$，是一个微妙的问题，将在第 10 节讨论（这正是选择记号 $\mathcal{A}_1,\mathcal{A}_2$ 而不选 $\mathcal{A}_+,\mathcal{A}_-$ 的原因）。

$\mathbb{Z}_2$ 奇点的轨迹维数为 $2g-2$。在选取谱曲线 $z^2+\det\varphi = 0$（其中 $\det\varphi$ 只有二重零点）时有 $g-1$ 个参数，而在选取适当的线丛 $\mathcal{L}'\to D'$ 时又有 $g-1$ 个参数。

% ==== 原文 PDF 第 65 页（印刷页 65）====

\subsubsection{转移}

现在让我们从 $SO_3$ 局部系统的观点来看这一结果。若一个 $SO_3$ 局部系统的结构群约化到 $O_2$、但不约化到其真子群，则它应当表示 $\mathcal{M}_H$ 的一个 $\mathbb{Z}_2$ 轨形奇点。$O_2$ Higgs 丛的模空间的维数为 $2g-2$，与上述轨形奇点族的维数一致。

更具体地说，表示奇点的那些 Higgs 丛正是 $O_2$ Higgs 丛。$\mathcal{M}_H(O_2)$ 就是 $\mathcal{M}_H(O_2^*)$ 对 $Q = H^1(C,\mathbb{Z}_2)$ 的商。故其 Hitchin 纤维化容易理解：底与 $O_2^*$ 的相同，而纤维是 Prym 对 $Q$ 的商。这与我们刚才描述的 $\mathcal{M}_H(SO_3)$ 的奇异轨迹相同。

与式 (5.15) 周围的论证相同的论证表明，即使在指定了非分歧二重覆盖 $D'\to C$ 之后，$\mathcal{M}_H(O_2)$ 仍有两个分支。事实上，从 $O_2$ 局部系统 $U$ 的观点看，非分歧二重覆盖的选择指定了 Stieffel--Whitney 类 $w_1(U)$，而对给定 $w_1(U)$ 的选择，两个分支由 $w_2(U)$ 的值区分。从 $O_2$ 局部系统到 $SO_3$ 局部系统的映射 $U\to U\oplus\det U$ 消去 $w_1$ 而保持 $w_2$ 不变。于是 $\mathcal{M}_H(O_2)$ 中具有给定 $w_2$ 值的所有分支都作为 $\mathbb{Z}_2$ 轨形奇点出现在 $\mathcal{M}_H(SO_3)$ 中以同一 $w_2$ 值标记的分支中。

$\mathcal{M}_H(O_2)$ 作为 $\mathcal{M}_H(SO_3)$ 的奇异轨迹出现，这一事实对 $O_2$ 与 $SO_3$ 这一对群并不特殊。对任何约化 Lie 群 ${}^{L}G$ 与约化子群 ${}^{L}H$，都有自然嵌入 $\mathcal{M}_H({}^{L}H)\subset\mathcal{M}_H({}^{L}G)$。若 ${}^{L}H$ 在 ${}^{L}G$ 中的中心化子非平凡，则 $\mathcal{M}_H({}^{L}H)$ 将是一条奇异轨迹。这一嵌入导致从 $\mathcal{M}_H({}^{L}H)$ 上 B-膜范畴到 $\mathcal{M}_H({}^{L}G)$ 上 B-膜范畴的自然函子（直接像），而后者必将给出从 $\mathcal{M}_H(H)$ 上 A-膜范畴到 $\mathcal{M}_H(G)$ 上 A-膜范畴的函子，如下图所示：
\[
\begin{array}{ccc}
\mathcal{M}_H({}^{L}G)\text{ 上的 }\mathcal{B}\text{-膜} & \xrightarrow{\ \sim\ } & \mathcal{M}_H(G)\text{ 上的 }\mathcal{A}\text{-膜}\\[10pt]
\uparrow & & \uparrow\ \text{transfer}\\[10pt]
\mathcal{M}_H({}^{L}H)\text{ 上的 }\mathcal{B}\text{-膜} & \xrightarrow{\ \sim\ } & \mathcal{M}_H(H)\text{ 上的 }\mathcal{A}\text{-膜}
\end{array}
\]
% [Tutor Note: 推导补足。图中两个水平箭头来自 Hitchin 纤维化的镜像对称（[KW] 的范畴等价），而右侧竖直箭头「transfer」来自把 B-膜沿 $\mathcal{M}_H({}^{L}H)\hookrightarrow\mathcal{M}_H({}^{L}G)$ 作直接像。作者说这个箭头是「必将给出」的，其逻辑是：左竖直箭头（B-膜到 B-膜）由几何直接给出，两个水平箭头是等价，于是复合给出右竖直箭头（A-膜到 A-膜）。注意图的方向：底部（小群 $H$）到顶部（大群 $G$），与经典情形中从 $H(\mathbb{A})$ 的 L-包到 $G(\mathbb{A})$ 的 L-包的转移方向一致。 来源：原书 p.65 上下文（A）。]

这与 Langlands 纲领中所谓的转移或函子性原理密切相关（见 [L1, Art1]）。在经典情形（第 7 节有更详细的讨论），这意味着：对偶群的任何同态 ${}^{L}H\to{}^{L}G$ 都期望给出一个映射（「转移」），从 $H(\mathbb{A})$ 的不可约自守表示的等价类集合（更确切地说，其 L-包）到 $G(\mathbb{A})$ 的相应集合。在几何情形，自守表示被替换为 $\Bun_G$ 上的 $\CD$-模，或 $\mathcal{M}_H(G)$ 上的 A-膜，而转移成为联系于 $H$ 与 $G$ 的适当范畴之间的函子，如上图所示。该函子应当与 Hecke/'t Hooft 算子（下一节讨论）在两个范畴上的作用相容。

Hitchin 纤维化的镜像对称为构造这样的函子提供了自然的框架。用物理的语言，人们或许希望通过引入超对称域壁来研究这一情形：域壁的一侧是规范群为 ${}^{L}H_c$（${}^{L}H$ 的紧实形式）的 $N=4$ 超 Yang--Mills 理论，另一侧是同样的
% ==== 原文 PDF 第 66 页（印刷页 66）====
规范群为 ${}^{L}G_c$ 的理论。关于这类域壁的弦论构造的一些例子，见 [Wi1]。

在我们的例子中 ${}^{L}H = O_2$ 而 ${}^{L}G = SO_3$。在对偶一侧我们有 $G = SL_2$，而 $H(F)$ 是 $SL_2(F)$ 中的一个扭环面，其中 $F$ 是 $C$ 上的有理函数域，其定义如下。考虑 $w_2=0$ 的 $O_2$-Higgs 丛的模空间 $\mathcal{M}_H(O_2,0)$，它嵌入 $\mathcal{M}_H(SO_3,0)$。它的分支由非分歧二重覆盖 $D'\to C$ 的集合 $J_2$ 参数化。对每个对应于某个非分歧覆盖的 $\psi\in J_2^\times$，我们定义 $H_\psi(F)$ 为 $D'$ 上非零有理函数 $f$ 构成的群，使得 $\tau(f) = f^{-1}$，其中 $\tau$ 是覆盖的对合。它自然地实现为 $SL_2(F)$ 的子群。上图所示转移函子则构造如下。

对应于 $\psi\in J_2^\times$ 的 $\mathcal{M}_H(O_2,0)$ 的每个分支 $\mathcal{M}_H(O_2,0)_\psi$ 是 Hitchin 底 $\boldsymbol{B}$ 中相应轨迹 $\boldsymbol{B}_\psi$ 上的环面纤维化。我们有 $\boldsymbol{B}_\psi = H^0(C,K\otimes\mathcal{L}_\psi)$，其中 $\mathcal{L}_\psi$ 是对应于 $\psi$ 的 $C$ 上的线丛，而映射 $\boldsymbol{B}_\psi\to\boldsymbol{B} = H^0(C,K^2)$ 由 $\omega\mapsto\omega^2$ 给出。由于我们希望避开约化到 $O_2$ 的真子群的局部系统，我们考虑 $\mathcal{M}_H(O_2,0)_\psi$ 中零纤维的补集 $\mathcal{M}_H(O_2,0)_\psi^\circ$。这是 $\mathcal{M}_H(SO_3,0)$ 中投影到 $\boldsymbol{B}_\psi^\circ = (\boldsymbol{B}_\psi\setminus 0)\subset\boldsymbol{B}$ 的子簇。这些子簇的并恰是 $\mathcal{M}_H(SO_3,0)$ 的「椭圆内窥轨迹」（见 [N1] 与第 9 节）。对每个点 $b\in\boldsymbol{B}_\psi^\circ$，Hitchin 纤维 $\boldsymbol{F}_b(O_2)_\psi\subset\mathcal{M}_H(O_2,0)_\psi$ 与对偶模空间 $\mathcal{M}_H(SL_2)$ 中相应奇异 Hitchin 纤维 $\boldsymbol{F}_b$ 的两个分支 $\boldsymbol{F}_{b,1}$ 与 $\boldsymbol{F}_{b,2}$ 上各自的秩一酉局部系统的模空间等同。事实上，如 5.2.2 节所解释的，每个 $\boldsymbol{F}_{b,i}$ 同构于阿贝尔簇 $\boldsymbol{F}_b(O_2)_\psi^\vee$ 上的 $(\mathbb{CP}^1)^{2g-2}$-丛，而该阿贝尔簇对偶于 $\boldsymbol{F}_b(O_2)_\psi$。转移（在零 Hitchin 纤维之外）则通过环面纤维化 $\mathcal{M}_H(O_2,0)_\psi^\circ$ 的逐纤维 T-对偶实现。特别地，支撑在点 $E\in\boldsymbol{F}_b(O_2)_\psi\subset\mathcal{M}_H(O_2,0)_\psi^\circ$ 上的点层 B-膜给出 $\boldsymbol{F}_b\subset\mathcal{M}_H(SL_2)$ 上的一个磁本征膜，它是 $\boldsymbol{F}_b(O_2)_\psi^\vee$ 上相应的秩一酉局部系统到 $\boldsymbol{F}_{b,1}$ 与 $\boldsymbol{F}_{b,2}$ 的拉回之和。

% [背景补充: 「椭圆内窥」（elliptic endoscopy）指内窥局部系统所对应的谱曲线其正规化是非分歧二重覆盖、从而 Hitchin 纤维的奇点仅为轨形奇点的情形。这是 Ngô 理论中「椭圆点」处的内窥，也是本文能用镜像对称直接处理的情形；一般（非椭圆）内窥的几何对应（即原文 §9 预告的内容）仍在探索中。 译者按（背景补充，无外部依据）。]

% [Tutor Note: 物理图像。逐纤维 T-对偶把 $\boldsymbol{F}_b(O_2)_\psi$ 上的点（即一个点层 B-膜）映为对偶环面 $\boldsymbol{F}_b(O_2)_\psi^\vee$ 上的一个平坦酉局部系统（A-膜）；再沿两个分支 $\boldsymbol{F}_{b,1},\boldsymbol{F}_{b,2}$ 分别拉回并相加，就得到 $\boldsymbol{F}_b$ 上的 A-膜。称之为「磁」本征膜，是因为它对应于 't Hooft 线算子（磁）的本征膜，而与 Wilson 线算子（电）的本征膜相对；本文 §5.3 将说明 't Hooft 算子正是 Hecke 算子的物理体现。零点层（零 Hitchin 纤维）被排除，是因为那里 T-对偶退化。 来源：原书 p.66 上下文（A）。]

此外，$\mathcal{M}_H(O_2)_\psi^\circ$ 作为 $\mathbb{Z}_2$-轨形奇点的轨迹嵌入 $\mathcal{M}_H(SO_3)$。这导致支撑在该分支上的 B-膜范畴的「倍增」。在对偶一侧，这反映为对偶 Hitchin 纤维有两个分支 $\boldsymbol{F}_{b,1}$ 与 $\boldsymbol{F}_{b,2}$，同样导致相应 A-膜范畴的「倍增」。
% [存疑: 此处原文作 $\mathcal{M}_H(O_2)_\psi^\circ$（$O_2$ 后无 ",0"），而同一段前面写作 $\mathcal{M}_H(O_2,0)_\psi^\circ$；疑为同一对象的简写，照录原文。 来源：原书 p.66 英文原文；存疑待核（A）。]

更一般的例子将在第 9 节考虑。
